\documentclass[a4paper,landscape,headrule,footrule]{foils}

%%
%%% macros for Theories of Grammar
%%%
\usepackage{polyglossia}
\setdefaultlanguage{english}
%\setmainfont{TeX Gyre Pagella}


\newcommand{\logo}{~}
\newcommand{\header}[3]{%
\title{\vspace*{-2ex} \large JMORF/MORS --- Morpho-Syntax
\\[2ex] \Large  \emp{#2} \\ \emp{#3}}
\author{\blu{Francis Bond}   \\ 
\normalsize  \textbf{Palacký University}\\
\normalsize  \url{https://fcbond.github.io/}\\
\normalsize  \texttt{bond@ieee.org}}
\MyLogo{JMORF (2026)}
\renewcommand{\logo}{#2}
\hypersetup{
   pdfinfo={
     Author={Francis Bond},
     Title={#1: #2},
     Subject={JMORF/MORS: Morpho-Syntax},
     Keywords={Syntax, Morphology, HPSG, Unification, Constructions},
     License={CC BY 4.0}
   }
 }

\date{#1
  \\ Location: KB-1.11}
}
\usepackage{hyperref}
\hypersetup{
     colorlinks,
     linkcolor={blue!50!black},
     citecolor={red!50!black},
     urlcolor={blue!80!black}
   }



\usepackage{xcolor}
\usepackage{graphicx}
\newcommand{\blu}[1]{\textcolor{blue}{#1}}
\newcommand{\grn}[1]{\textcolor{green}{#1}}
\newcommand{\hide}[1]{\textcolor{white}{#1}}
\newcommand{\emp}[1]{\textcolor{red}{#1}}
\newcommand{\txx}[1]{\textbf{\textcolor{blue}{#1}}}
\newcommand{\lex}[1]{\textbf{\mtcitestyle{#1}}}


% \renewcommand{\labelitemi}{\textcolor{violet}{‣}}
% \renewcommand{\labelitemii}{\textcolor{purple}{‣}}

\newcommand{\subhead}[1]{\noindent\textbf{#1}\\[5mm]}

\newcommand{\Bad}{\emp{\raisebox{0.15ex}{\ensuremath{\mathbf{\otimes}}}}}
\newcommand{\bad}[1]{*\eng{#1}}

\newcommand{\com}[1]{\hfill (\emp{#1})}%

\usepackage{relsize,xspace}
\newcommand{\into}{\ensuremath{\rightarrow}\xspace}
\newcommand{\tot}{\ensuremath{\leftrightarrow}\xspace}
\usepackage{url}
\newcommand{\lurl}[1]{\MyLogo{\url{#1}}}

\usepackage{langsci-gb4e}
\let\eachwordone=\slshape
%%% langsci-gb4e defines \thexnumi as \@roman\@xsi{xnumi}, which only works
%%% inside an exe environment; some decks use xlisti on its own.  Restore
%%% the plain gb4e definition so that keeps working.
\makeatletter
\def\thexnumi{\@xsi{xnumi}}
\makeatother
\newcommand{\lx}[1]{\textbf{\mtciteform{#1}}}
\newcommand{\ix}{\ex\slshape}



\newcommand{\ent}{\ensuremath{\Rightarrow}\xspace}
\newcommand{\ngv}{\ensuremath{\not\Rightarrow}\xspace}
%\usepackage{times}
%\usepackage{nttfoilhead}
\newcommand{\myslide}[1]{\foilhead[-25mm]{\raisebox{12mm}[0mm]{\emp{#1}}}\MyLogo{\logo}}
\newcommand{\myslider}[1]{\rotatefoilhead[-25mm]{\raisebox{12mm}[0mm]{\emp{#1}}}}
%\newcommand{\myslider}[1]{\rotatefoilhead{\raisebox{-8mm}{\emp{#1}}}}

\newcommand{\section}[1]{\myslide{}{\begin{center}\Huge \emp{#1}\end{center}}}



\usepackage[lyons,j,e,k]{mtg2e}
%\renewcommand{\mtcitestyle}[1]{\textcolor{-red!75!green!50}{\textsl{#1}}}
\renewcommand{\mtcitestyle}[1]{\textcolor{teal}{\textsl{#1}}}
\newcommand{\cs}{\mtciteform}
\newcommand{\iz}[1]{\texttt{\textup{#1}}}
\newcommand{\gm}{\textsc}
\usepackage[normalem]{ulem}
\newcommand{\ul}{\uline}
\newcommand{\ull}{\uuline}
\newcommand{\wl}{\uwave}
\newcommand{\vs}{\ensuremath{\Leftrightarrow}~}
%%%
%%% Bibliography
%%%
\usepackage{natbib}
%\usepackage{url}
\usepackage{bibentry}



\usepackage[utf8]{inputenc}
%%% trees
\usepackage{multicol}
%%% pst-node: one example in lec-12-control draws a connector between
%%% two points with \rnode/\ncbar.  This used to come in with rtrees.
\usepackage{pst-node}
\usepackage{forest}
\useforestlibrary{linguistics}
\forestapplylibrarydefaults{linguistics}

%%% \makebox's optional arguments are read as AVM delimiters by
%%% langsci-avm, so hide them behind a macro with no visible brackets.
\newcommand{\lbox}[1]{\makebox[0mm][l]{#1}}

%%% Tree styles.
%%%   words  -- trees whose leaves are actual words (italicised, aligned)
%%%   cat    -- trees whose leaves are categories or feature structures
\forestset{
  words/.style={for tree={s sep=7mm, l sep=5mm,
                          if n children=0{font=\itshape, tier=terminal}{}}},
  cat/.style={for tree={s sep=7mm, l sep=5mm}},
}
\newcommand{\wtree}[1]{\begin{forest} words #1 \end{forest}}
\newcommand{\ctree}[1]{\begin{forest} cat #1 \end{forest}}
%%% the rule being applied at this step of a derivation
\newcommand{\step}[1]{\begin{center}\emp{#1}\end{center}}

%%% AVMs (Language Science Press)
\usepackage{langsci-avm}

%%% Macros the decks use, carried over from headx.tex
\newcommand{\blank}{\rule{3em}{1pt}\xspace}
\newcommand{\ft}[1]{\textsc{#1}}
\newcommand{\idn}{\mtciteform}    %%% Indonesian
\newcommand{\prd}[1]{\textbf{#1}}
\newcommand{\typ}[1]{\textit{#1}}
%%% \val: provide before renewing, in case a package defines it.
\providecommand{\val}[1]{}
\renewcommand{\val}[1]{\textit{#1}}
%%% Manning's \avmbox, still used in tree labels and in maths where a
%%% bare langsci \tag would not work.  \avm{\tag{...}} gives the same
%%% boxed tag as the ones inside feature structures.
\newcommand{\avmbox}[1]{\mbox{\avm{\tag{#1}}}}

%%% Ring a value in a dotted oval.  These used pstricks \ovalnode,
%%% which came in with rtrees; tikz (via forest) does the same job.
\newcommand{\OV}[1]{\tikz[baseline]{\node[draw,dotted,red,ellipse,
          inner sep=1pt,outer sep=0pt]{#1};}}
\newcommand{\OVB}[1]{\tikz[baseline]{\node[draw,dotted,blue,ellipse,
          inner sep=1pt,outer sep=0pt]{#1};}}

%%% From CSLI book
\newcommand{\mc}{\multicolumn}
\newcommand{\HD}{\textbf{H}\xspace}
\newcommand{\el}{< >}       %%% empty list; langsci-avm needs the space

\newcommand{\A}{\noindent\textbf{A}: }
\newcommand{\Q}{\noindent\textbf{Q}: }
%\newcommand{\C}{\noindent\textbf{C}: }



\begin{document}

\header{Lecture 7}{The Structure of the Lexicon}{}
\maketitle


\myslide{Overview}
\begin{itemize}
\item Lexical types: the lexeme hierarchy and default inheritance
\item Lexical rules: inflectional, derivational and post-inflectional
\item Problem: arguments in Japanese
\end{itemize}

%%%%%%%%%%%%%%%%%%%%%%%%%%%%%%%%%%%%%%%%%%%%%%%%%%%%%%%%%%%%%%%%%%%%%%%%
%%% Converted from Emily Bender's slides (ch08a-uw.pdf), with permission
%%%%%%%%%%%%%%%%%%%%%%%%%%%%%%%%%%%%%%%%%%%%%%%%%%%%%%%%%%%%%%%%%%%%%%%%
\section{Lexical Types}

%%% p2
\myslide{Overview}
\begin{itemize}
\item Motivation for lexical hierarchy
\item Default inheritance
\item Tour of the lexeme hierarchy
\item The Case Constraint
\item \val{pos} vs.\ \val{lexeme}
\end{itemize}

%%% p3
\myslide{Motivation}
\begin{itemize}
\item We've streamlined our grammar rules\ldots
  \begin{itemize}
  \item \ldots by stating some constraints as general principles
  \item \ldots and locating lots of information in the lexicon.
  \item Our lexical entries currently stipulate a lot of
    information that is common across many entries and
    should be stated only once.
  \end{itemize}
\item Examples?
\item Ideally, particular lexical entries need only give
  phonological form, the semantic contribution,
  and any constraints truly idiosyncratic to the
  lexical entry.
\end{itemize}

%%% p4
\myslide{Lexemes and Words}
\begin{itemize}
\item \txx{Lexeme}: An abstract proto-word which gives rise
  to genuine words. We refer to lexemes by their
  `dictionary form', e.g.\ `the lexeme \eng{run}' or `the
  lexeme \eng{dog}'.
\item \txx{Word}: A particular pairing of form and meaning.
  \eng{Running} and \eng{ran} are different words
\end{itemize}

%%% p5
\myslide{Lexical Types \& Lexical Rules}
\begin{itemize}
\item Lexemes capture the similarities among \lex{run, runs,
    running,} and \lex{run}.
\item The lexical type hierarchy captures the similarities among
  \lex{run, sleep,} and \lex{laugh}, among those and other verbs like
  \lex{devour} and \lex{hand}, and among those and other words like
  \lex{book}.
  \begin{itemize}
  \item[Q] What do \lex{devour} and \lex{book} have in common?
  \item[A] The SHAC
  \end{itemize}
\item Lexical rules capture the similarities among \lex{runs, sleeps,
    devours, hands,}\ldots
\end{itemize}

%%% p6
\myslide{Default Inheritance}
\begin{itemize}
\item[Q] Why do we have default inheritance?
\item[A] Generalizations with exceptions are common:
  \begin{itemize}
  \item Most nouns in English aren't marked for \ft{case}, but
    pronouns are.
  \item Most verbs in English only distinguish two agreement
    categories (\val{3sing} and \val{non-3sing}), but \lex{be} distinguishes
    more.
  \item Most prepositions in English are transitive, but \lex{here} and
    \lex{there} are intransitive.
  \item Most nominal words in English are 3rd person, but some
    (all of them pronouns) are 1st or 2nd person.
  \item Most proper nouns in English are singular, but some
    (mountain range names, sports team names) are plural.
  \end{itemize}
\end{itemize}

%%% p7
\myslide{Default Inheritance, Technicalities}
\vspace*{1em}
\begin{tabular}{@{}p{0.28\textwidth}p{0.28\textwidth}p{0.32\textwidth}@{}}
\raggedright If a type says \avm{[ arg-st & / < NP > ]}, &
\raggedright and one of its subtypes says \avm{[ arg-st & < > ]}, &
\raggedright then the \ft{arg-st} value of instances of the subtype is
\avm{< >}.\tabularnewline[3em]
\raggedright If a type says \avm{[ arg-st & < NP > ]}, &
\raggedright and one of its subtypes says \avm{[ arg-st & < > ]}, &
\raggedright then this subtype can have no instances, since they would have to
satisfy contradictory constraints.\tabularnewline
\end{tabular}

%%% p8
\myslide{Default Inheritance, More Technicalities}
\begin{itemize}
\item If a type says \avm{[ mod & / < S > ]}, and one of its subtypes says
  \avm{[ mod & < [ spr & < NP > ] > ]}, then the \ft{arg-st} value of
  instances of the subtype is what?
\end{itemize}
\begin{center}
  \avm{[ mod & < [ head & / \type{verb} \\
                   spr & < NP > \\
                   comps & / < > ] > ]}
\end{center}
\begin{itemize}
\item That is, default constraints are `pushed down'
\end{itemize}

%%% p9
\myslide{Question on Default Inheritance}
\begin{itemize}
\item[Q] Can a grammar rule override a default
  constraint on a word?
\item[A] No. Defaults are all `cached out' in the
  lexicon.
\item Words as used to build sentences have only
  inviolable constraints.
\end{itemize}

%%% The lexeme hierarchy is shown repeatedly, with an arrow pointing
%%% at the type under discussion.  #1: comma-separated names of nodes
%%% to colour red; #2: TikZ code drawn after the tree (the arrow).
\makeatletter
\newcommand{\lxhl}[2]{\ifcsname lexhl@#1\endcsname\textcolor{red}{#2}\else #2\fi}
\newcommand{\lexhier}[2]{%
\begingroup
\@for\lx@n:=#1\do{\expandafter\def\csname lexhl@\lx@n\endcsname{}}%
\begin{center}
\scalebox{0.7}{\begin{forest}
cat, for tree={font=\itshape, s sep=3mm}
[{\lxhl{synsem}{synsem}\\\textnormal{\avm{[ \textsc{syn, sem} ]}}}, align=center, name=synsem, s sep=4mm
  [{lexeme\\\textnormal{\avm{[ arg-st ]}}}, align=center, name=lexeme
    [{\lxhl{infl}{infl-lxm}}, name=infl
      [{\lxhl{verb}{verb-lxm}}, l*=3, name=verb
        [{siv-lxm}, name=siv]
        [{piv-lxm}, name=piv]
        [{tv-lxm}, name=tv
          [{stv-lxm}, name=stv]
          [{dtv-lxm}, name=dtv]
          [{ptv-lxm}, name=ptv]]]
      [{\lxhl{cn}{cn-lxm}}, l*=3, name=cn
        [{\lxhl{cntn}{cntn-lxm}}, name=cntn]
        [{\lxhl{massn}{massn-lxm}}, name=massn]]]
    [{const-lxm}, name=const
      [{adj-lxm}, l*=2, name=adj]
      [{conj-lxm}, l*=2, name=conj]
      [{det-lxm}, l*=2, name=det]
      [{predp-lxm}, l*=2, name=predp]
      [{argmkp-lxm}, l*=2, name=argmkp]
      [{\lxhl{pn}{pn-lxm}}, l*=1.4, name=pn]
      [{\lxhl{pron}{pron-lxm}}, l*=1.4, name=pron]]]
  [{expression}, name=expression, l*=0.8
    [{word\\\textnormal{\avm{[ arg-st ]}}}, align=center, l*=0.6, name=word]
    [{\lxhl{phrase}{phrase}}, l*=0.6, name=phrase]]]
#2
\end{forest}}
\end{center}
\endgroup}
\makeatother
%%% the orange arrow
\tikzset{hlarrow/.style={overlay, ->, >=stealth, orange, line width=2mm}}

%%% p10
\myslide{Our Lexeme Hierarchy}
\lexhier{}{}

%%% p11
\myslide{Functions of Types}
\begin{itemize}
\item Stating what features are appropriate for
  what categories
\item Stating generalizations
  \begin{itemize}
  \item Constraints that apply to (almost) all instances
  \item Generalizations about selection --- where
    instances of that type can appear
  \end{itemize}
\end{itemize}

%%% p12
\myslide{Every \val{synsem} has the Features SYN and SEM}
\lexhier{synsem}{\draw[hlarrow] ([xshift=-25mm]synsem.west) -- ([xshift=-2mm]synsem.west);}

%%% p13
\myslide{No ARG-ST on \val{phrase}}
\lexhier{phrase}{\draw[hlarrow] ([xshift=20mm,yshift=-12mm]phrase.south) -- ([yshift=-2mm]phrase.south);}

%%% p14
\myslide{A Constraint on \val{infl-lxm}: the SHAC}
\lexhier{infl}{\draw[hlarrow] ([xshift=-25mm]infl.west) -- ([xshift=-2mm]infl.west);}

%%% p15
\myslide{A Constraint on \val{infl-lxm}: the SHAC}
\vspace*{2em}
\begin{center}
\val{infl-lxm}: \avm{[ syn & [ val & [ spr & < [ agr & \1 ] > ] \\
                               head & [ agr & \1 ] ] ]}
\end{center}

%%% p16
\myslide{Constraints on \val{cn-lxm}}
\lexhier{cn}{\draw[hlarrow] ([xshift=25mm]cn.east) -- ([xshift=2mm]cn.east);}

%%% p17
\myslide{Constraints on \val{cn-lxm}}
\begin{center}
\val{cn-lxm}: \avm{[ syn & [ head & [\type*{noun}
                                     agr & [ per & 3rd ] ] \\
                             val & [ spr & < [ head & \type{det} \\
                                               index & \textit{i} ] > ] ] \\
                     sem & [ mode & / ref \\
                             index & \textit{i} ] \\
                     arg-st & < X > $\oplus$ / < > ]}
\end{center}

%%% p18
\myslide{Formally Distinguishing Count vs.\ Mass Nouns}
\lexhier{cntn,massn}{%
  \draw[hlarrow] ($(cntn.south)!0.5!(massn.south)+(11mm,-12mm)$) -- ([yshift=-2mm]cntn.south);
  \draw[hlarrow] ($(cntn.south)!0.5!(massn.south)+(13mm,-12mm)$) -- ([yshift=-2mm]massn.south);}

%%% p19
\myslide{Formally Distinguishing Count vs.\ Mass Nouns}
\vspace*{2em}
\begin{center}
\begin{tabular}{r@{\ }l}
\val{cntn-lxm}: & \avm{[ syn & [ val & [ spr & < [ count & + ] > ] ] ]}\\[3ex]
\val{massn-lxm}: & \avm{[ syn & [ val & [ spr & < [ count & $-$ ] > ] ] ]}
\end{tabular}
\end{center}

%%% p20
\myslide{Constraints on \val{verb-lxm}}
\lexhier{verb}{\draw[hlarrow] ([xshift=-15mm,yshift=20mm]verb.north west) -- (verb.north west);}

%%% p21
\myslide{Constraints on \val{verb-lxm}}
\vspace*{2em}
\begin{center}
\val{verb-lxm}: \avm{[ syn & [ head & \type{verb} ] \\
                       sem & [ mode & prop ] \\
                       arg-st & / < NP, \ldots > ]}
\end{center}

%%% p22
\myslide{Subtypes of \val{verb-lxm}}
\begin{center}
\scalebox{0.8}{\begin{forest}
cat, for tree={font=\itshape, s sep=3mm, l sep=3mm}
[{verb-lxm}
  [{siv-lxm}]
  [{piv-lxm}]
  [{tv-lxm}
    [{stv-lxm}]
    [{dtv-lxm}]
    [{ptv-lxm}]]]
\end{forest}}
\end{center}
\vspace*{-1ex}
\begin{itemize}
\item \val{verb-lxm}: \avm{[ arg-st & / < NP, \ldots > ]}
  \begin{itemize}
  \item \val{siv-lxm}: \avm{[ arg-st & / < NP > ]}
  \item \val{piv-lxm}: \avm{[ arg-st & / < NP, PP > ]}
  \item \val{tv-lxm}: \avm{[ arg-st & / < NP, NP, \ldots > ]}
    \begin{itemize}
    \item \val{stv-lxm}: \avm{[ arg-st & / < NP, NP > ]}
    \item \val{dtv-lxm}: \avm{[ arg-st & / < NP, NP, NP > ]}
    \item \val{ptv-lxm}: \avm{[ arg-st & / < NP, NP, PP > ]}
    \end{itemize}
  \end{itemize}
\end{itemize}

%%% p23
\myslide{Proper Nouns and Pronouns}
\lexhier{pn,pron}{%
  \draw[hlarrow] ([xshift=4mm,yshift=-18mm]pn.south) -- ([yshift=-2mm]pn.south);
  \draw[hlarrow] ([xshift=12mm,yshift=-18mm]pron.south) -- ([yshift=-2mm]pron.south);}

%%% p24
\myslide{Proper Nouns and Pronouns}
\begin{center}
\begin{tabular}{r@{\ }l}
\val{pn-lxm}: &
\avm{[ syn & [ head & [\type*{noun}
                       agr & [ per & 3rd \\
                               num & / sg ] ] ] \\
       sem & [ mode & ref ] \\
       arg-st & / < > ]}\\[4ex]
\val{pron-lxm}: &
\avm{[ syn & [ head & \type{noun} ] \\
       sem & [ mode & / ref ] \\
       arg-st & < > ]}
\end{tabular}
\end{center}

%%% p25
%%% a green tick (no amssymb \checkmark here)
\newcommand{\tick}{\tikz\draw[green!60!black, line width=2pt] (0,0.4em) -- (0.3em,0) -- (0.8em,0.9em);}
\myslide{The Case Constraint}
An outranked NP is \avm{[ case & acc ]}.
\begin{itemize}
\item object of verb \hfill \tick\hspace*{3.2em}
\item second object of verb \hfill \tick\hspace*{3.2em}
\item object of argument-marking preposition \hfill \tick\hspace*{3.2em}
\item object of predicational preposition \hfill (\tick)\hspace*{2.8em}
\end{itemize}

%%% p26
\myslide{The Case Constraint, Continued}
An outranked NP is \avm{[ case & acc ]}.
\begin{itemize}
\item Subjects of verbs
  \begin{itemize}
  \item Should we add a clause to cover nominative subjects?
    \begin{itemize}
    \item No.\\
      \eng{We expect them to leave.} (Chapter 12)
    \item Lexical rules for finite verbs will handle nominative subjects.
    \end{itemize}
  \end{itemize}
\item Any other instances of case marking in English?
\item Does it apply to case systems in other languages?\\
  No: The Case Constraint is an English-specific constraint.
\end{itemize}

%%% p27
\myslide{Apparent Redundancy}
\begin{itemize}
\item Why do we need both the \val{pos} subhierarchy
  and lexeme types?
\item \val{pos}:
  \begin{itemize}
  \item Applies to words and phrases; models
    relationship between them
  \item Constrains which features are appropriate
    (no \ft{aux} on \val{noun})
  \end{itemize}
\item \val{lexeme}:
  \begin{itemize}
  \item Generalizations about combinations of
    constraints
  \end{itemize}
\end{itemize}

%%% p28
\myslide{Lexical Types \& Lexical Rules}
\begin{itemize}
\item Lexemes capture the similarities among \lex{run, runs,
    running,} and \lex{run}.
\item The lexical type hierarchy captures the similarities among
  \lex{run, sleep,} and \lex{laugh}, among those and other verbs like
  \lex{devour} and \lex{hand}, and among those and other words like
  \lex{book}.
\item Lexical rules capture the similarities among \lex{runs, sleeps,
    devours, hands,}\ldots
\end{itemize}

%%% p29
\myslide{Overview}
\begin{itemize}
\item Motivation for lexical hierarchy
\item Default inheritance
\item Tour of the lexeme hierarchy
\item The Case Constraint
\item \val{pos} vs.\ \val{lexeme}
\end{itemize}

%%%%%%%%%%%%%%%%%%%%%%%%%%%%%%%%%%%%%%%%%%%%%%%%%%%%%%%%%%%%%%%%%%%%%%%%
%%% Converted from Emily Bender's slides (ch08b-uw.pdf), with permission
%%%%%%%%%%%%%%%%%%%%%%%%%%%%%%%%%%%%%%%%%%%%%%%%%%%%%%%%%%%%%%%%%%%%%%%%
\section{Lexical Rules}

%%% Grey box over a part of an AVM that is not yet revealed (her
%%% build-up slides cover parts of the AVM with grey rectangles).
\providecommand{\gry}[1]{{\setlength{\fboxsep}{0pt}\colorbox{lightgray}{\phantom{#1}}}}
%%% Boxes holding hidden AVM parts that contain & (these cannot be
%%% passed through \gry inside an \avm directly).
\ifdefined\gryA\else\newsavebox{\gryA}\fi
\ifdefined\gryB\else\newsavebox{\gryB}\fi
\ifdefined\gryC\else\newsavebox{\gryC}\fi

\myslide{Overview}
\begin{itemize}
\item How lexical rules fit in
\item Three types of lexical rules, constraints
\item Example: Plural noun lexical rule
\item Advice on writing lexical rules
\item Constant lexemes
\item \ft{arg-st} \& ARP
\item The feature \ft{form}
\end{itemize}

\myslide{Lexical Types \& Lexical Rules}
\begin{itemize}
\item Lexemes capture the similarities among
  \eng{run}, \eng{runs}, \eng{running}, and \eng{ran}
\item The lexical type hierarchy captures the
  similarities among \eng{run}, \eng{sleep}, and \eng{laugh},
  among those and other verbs like \eng{devour}
  and \eng{hand}, and among those and other words
  like \eng{book}.
\item Lexical rules capture the similarities among
  \eng{runs}, \eng{sleeps}, \eng{devours}, \eng{hands}, \ldots
\end{itemize}

\myslide{Parsimony \& Plausibility}
\begin{itemize}
\item Lexical rules capture \textbf{productive}
  generalizations.
\item There may be some `precompiling'
  going on as well.
\end{itemize}

\myslide{Three Kinds of Lexical Rules}
\begin{itemize}
\item \txx{Inflectional}: \val{lexeme} to \val{word}
  \\[1ex] Examples?
\item \txx{Derivational}: \val{lexeme} to \val{lexeme}
  \\[1ex] Examples?
\item \txx{Post-Inflectional}: \val{word} to \val{word}
  \\ (Chapters 11, 13, 14)
\end{itemize}

\myslide{Three Subtypes of \textit{l-rule}}
\begin{center}
\begin{small}
  \begin{forest}
    cat, for tree={font=\itshape, l sep=3mm}
    [{l-rule} [{i-rule}] [{d-rule}] [{pi-rule}]]
  \end{forest}

  \medskip
  \val{l-rule} :
  \avm{[ input & \type{l-sequence} < X , [ sem & / \2 ] > \\
         output & \type{l-sequence} < Y , [ sem & / \2 ] > ]}

  \medskip
  \val{i-rule} :
  \avm{[ input & < X , [\type*{lexeme}
                         syn & \3 \\
                         arg-st & \tag{A} ] > \\
         output & < Y , [\type*{word}
                          syn & \3 \\
                          arg-st & \tag{A} ] > ]}
  \qquad
  \val{d-rule} :
  \avm{[ input & < X , [\type*{lexeme}
                         syn & / \3 ] > \\
         output & < Y , [\type*{lexeme}
                          syn & / \3 ] > ]}
\end{small}
\end{center}

\myslide{Plural Noun LR}
\begin{center}
  \avm{[\type*{i-rule}
    input & < \1 , \type{cntn-lxm} > \\
    output & < F$_{NPL}$(\1) , [\type*{word}
                syn & [ head & [ agr & [ num & pl ] ] ] ] > ]}
\end{center}

\myslide{Plural Noun LR with Inherited Constraints}
\begin{center}
\begin{small}
\sbox{\gryA}{\avm{[ head & [\type*{noun} agr & \4 [ per & 3rd ] ] \\
                  val & [ spr & < DP[ count & + \\ agr & \4 ] > ] ]}}
\sbox{\gryB}{\avm{[ mode & / ref ]}}
\sbox{\gryC}{\avm{[ spr & \tag{B} \\ comps & \tag{C} ]}}
  \avm{[\type*{i-rule}
    input & < \1 , [\type*{cntn-lxm}
       \gry{syn} & \gry{\avm{\3}} \gry{\usebox{\gryA}} \\
       \gry{sem} & \gry{\avm{\2}} \gry{\usebox{\gryB}} \\
       \gry{arg-st} & \gry{\avm{\tag{B} $\oplus$ \tag{C}}} ] > \\
    output & < F$_{NPL}$(\1) , [\type*{word}
       \gry{syn} & \gry{\avm{\3}} [ head & [ agr & [ num & pl ] ] \\
                  \gry{val} & \gry{\usebox{\gryC}} ] \\
       \gry{sem} & \gry{\avm{\2}} \\
       \gry{arg-st} & \gry{\avm{\tag{B} $\oplus$ \tag{C}}} ] > ]}
\end{small}
\end{center}

\myslide{Plural Noun LR with Inherited Constraints}
\begin{center}
\begin{small}
\sbox{\gryA}{\avm{[ head & [\type*{noun} agr & \4 [ per & 3rd ] ] \\
                  val & [ spr & < DP[ count & + \\ agr & \4 ] > ] ]}}
\sbox{\gryB}{\avm{[ mode & / ref ]}}
\sbox{\gryC}{\avm{[ spr & \tag{B} \\ comps & \tag{C} ]}}
  \avm{[\type*{i-rule}
    input & < \1 , [\type*{cntn-lxm}
       \gry{syn} & \gry{\avm{\3}} \gry{\usebox{\gryA}} \\
       sem & \avm{\2} \gry{\usebox{\gryB}} \\
       \gry{arg-st} & \gry{\avm{\tag{B} $\oplus$ \tag{C}}} ] > \\
    output & < F$_{NPL}$(\1) , [\type*{word}
       \gry{syn} & \gry{\avm{\3}} [ head & [ agr & [ num & pl ] ] \\
                  \gry{val} & \gry{\usebox{\gryC}} ] \\
       sem & \2 \\
       \gry{arg-st} & \gry{\avm{\tag{B} $\oplus$ \tag{C}}} ] > ]}
\end{small}
\end{center}

\myslide{Plural Noun LR with Inherited Constraints}
\begin{center}
\begin{small}
\sbox{\gryA}{\avm{[ head & [\type*{noun} agr & \4 [ per & 3rd ] ] \\
                  val & [ spr & < DP[ count & + \\ agr & \4 ] > ] ]}}
\sbox{\gryB}{\avm{[ mode & / ref ]}}
\sbox{\gryC}{\avm{[ spr & \tag{B} \\ comps & \tag{C} ]}}
  \avm{[\type*{i-rule}
    input & < \1 , [\type*{cntn-lxm}
       \gry{syn} & \gry{\avm{\3}} \gry{\usebox{\gryA}} \\
       sem & \avm{\2} \gry{\usebox{\gryB}} \\
       arg-st & \tag{B} $\oplus$ \tag{C} ] > \\
    output & < F$_{NPL}$(\1) , [\type*{word}
       \gry{syn} & \gry{\avm{\3}} [ head & [ agr & [ num & pl ] ] \\
                  \gry{val} & \gry{\usebox{\gryC}} ] \\
       sem & \2 \\
       arg-st & \tag{B} $\oplus$ \tag{C} ] > ]}
\end{small}
\end{center}

\myslide{Plural Noun LR with Inherited Constraints}
\begin{center}
\begin{small}
\sbox{\gryA}{\avm{[ head & [\type*{noun} agr & \4 [ per & 3rd ] ] \\
                  val & [ spr & < DP[ count & + \\ agr & \4 ] > ] ]}}
\sbox{\gryB}{\avm{[ mode & / ref ]}}
\sbox{\gryC}{\avm{[ spr & \tag{B} \\ comps & \tag{C} ]}}
  \avm{[\type*{i-rule}
    input & < \1 , [\type*{cntn-lxm}
       syn & \avm{\3} \gry{\usebox{\gryA}} \\
       sem & \avm{\2} \gry{\usebox{\gryB}} \\
       arg-st & \tag{B} $\oplus$ \tag{C} ] > \\
    output & < F$_{NPL}$(\1) , [\type*{word}
       syn & \avm{\3} [ head & [ agr & [ num & pl ] ] \\
                  \gry{val} & \gry{\usebox{\gryC}} ] \\
       sem & \2 \\
       arg-st & \tag{B} $\oplus$ \tag{C} ] > ]}
\end{small}
\end{center}

\myslide{Plural Noun LR with Inherited Constraints}
\begin{center}
\begin{small}
\sbox{\gryC}{\avm{[ spr & \tag{B} \\ comps & \tag{C} ]}}
  \avm{[\type*{i-rule}
    input & < \1 , [\type*{cntn-lxm}
       syn & \avm{\3} [ head & [\type*{noun} agr & \4 [ per & 3rd ] ] \\
                  val & [ spr & < DP[ count & + \\ agr & \4 ] > ] ] \\
       sem & \avm{\2} [ mode & / ref ] \\
       arg-st & \tag{B} $\oplus$ \tag{C} ] > \\
    output & < F$_{NPL}$(\1) , [\type*{word}
       syn & \avm{\3} [ head & [ agr & [ num & pl ] ] \\
                  \gry{val} & \gry{\usebox{\gryC}} ] \\
       sem & \2 \\
       arg-st & \tag{B} $\oplus$ \tag{C} ] > ]}
\end{small}
\end{center}

\myslide{Plural Noun LR with Inherited Constraints}
\begin{center}
\begin{small}
  \avm{[\type*{i-rule}
    input & < \1 , [\type*{cntn-lxm}
       syn & \avm{\3} [ head & [\type*{noun} agr & \4 [ per & 3rd ] ] \\
                  val & [ spr & < DP[ count & + \\ agr & \4 ] > ] ] \\
       sem & \avm{\2} [ mode & / ref ] \\
       arg-st & \tag{B} $\oplus$ \tag{C} ] > \\
    output & < F$_{NPL}$(\1) , [\type*{word}
       syn & \avm{\3} [ head & [ agr & [ num & pl ] ] \\
                  val & [ spr & \OV{\avm{\tag{B}}} \\ comps & \OV{\avm{\tag{C}}} ] ] \\
       sem & \2 \\
       arg-st & \OV{\avm{\tag{B} $\oplus$ \tag{C}}} ] > ]}
\end{small}
\end{center}

\myslide{Practicalities --- Applying Lexical Rules}
\begin{itemize}
\item \ft{input} is a family of lexical sequences.
\item \ft{output} is another family of lexical sequences.
  \begin{itemize}
  \item \ldots usually a smaller family
  \item \ldots usually a disjoint one
  \end{itemize}
\item The only differences between the families are those
  stipulated in the rule (or the rule's type).
\item Similarities are handled by the constraints on \val{l-rule}
  and its subtypes.
\item If we've written the LRs correctly, nothing is left
  underconstrained.
\end{itemize}

\myslide{Example: Lexical Entry for \textit{cat}}
\begin{center}
  \avm{< \textnormal{cat} , [\type*{cntn-lxm}
    sem & [ index & $k$ \\
            restr & < [ reln & cat \\ inst & $k$ ] > ] ] >}
\end{center}

\myslide{Example: \textit{cat}, with Inheritance}
\begin{center}
\sbox{\gryA}{\avm{[\type*{noun} agr & \7 [ per & 3rd ] ]}}
\sbox{\gryB}{\avm{< DP[ count & + \\ agr & \7 ] >}}
  \avm{< \textnormal{cat} , [\type*{cntn-lxm}
    syn & [ \gry{head} & \gry{\usebox{\gryA}} \\
            val & [ \gry{spr} & \gry{\usebox{\gryB}} ] ] \\
    sem & [ \gry{mode} & \gry{ref} \\
            index & $k$ \\
            restr & < [ reln & cat \\ inst & $k$ ] > ] \\
    \gry{arg-st} & \gry{\avm{< X >}} ] >}
\end{center}

\myslide{Example: \textit{cat}, with Inheritance}
\begin{center}
\sbox{\gryA}{\avm{[\type*{noun} agr & \7 [ per & 3rd ] ]}}
  \avm{< \textnormal{cat} , [\type*{cntn-lxm}
    syn & [ \gry{head} & \gry{\usebox{\gryA}} \\
            val & [ spr & < \gry{DP}[ count & + \\ \gry{agr} & \gry{\avm{\7}} ] > ] ] \\
    sem & [ \gry{mode} & \gry{ref} \\
            index & $k$ \\
            restr & < [ reln & cat \\ inst & $k$ ] > ] \\
    \gry{arg-st} & \gry{\avm{< X >}} ] >}
\end{center}

\myslide{Example: \textit{cat}, with Inheritance}
\begin{center}
  \avm{< \textnormal{cat} , [\type*{cntn-lxm}
    syn & [ head & [\type*{noun} agr & \gry{\avm{\7}} [ per & 3rd ] ] \\
            val & [ spr & < DP[ count & + \\ \gry{agr} & \gry{\avm{\7}} ] > ] ] \\
    sem & [ mode & ref \\
            index & $k$ \\
            restr & < [ reln & cat \\ inst & $k$ ] > ] \\
    arg-st & < X > ] >}
\end{center}

\myslide{Example: \textit{cat}, with Inheritance}
\begin{center}
  \avm{< \textnormal{cat} , [\type*{cntn-lxm}
    syn & [ head & [\type*{noun} agr & \7 [ per & 3rd ] ] \\
            val & [ spr & < DP[ count & + \\ agr & \7 ] > ] ] \\
    sem & [ mode & ref \\
            index & $k$ \\
            restr & < [ reln & cat \\ inst & $k$ ] > ] \\
    arg-st & < X > ] >}
\end{center}

\myslide{Plural Noun LR}
\begin{center}
  \avm{[\type*{i-rule}
    input & < \1 , \type{cntn-lxm} > \\
    output & < F$_{NPL}$(\1) , [\type*{word}
                syn & [ head & [ agr & [ num & pl ] ] ] ] > ]}
\end{center}

\myslide{Licensing \textit{cats}}
\begin{center}
\scalebox{0.75}{%
  \avm{[\type*{i-rule}
    input & < \1\textnormal{cat} , [\type*{cntn-lxm}
       syn & \3 [ head & [\type*{noun} agr & \7 [ per & 3rd ] ] \\
                  val & [ spr & < DP[ count & + \\ agr & \7 ] > ] ] \\
       sem & \2 [ mode & ref \\
                  index & $k$ \\
                  restr & < [ reln & cat \\ inst & $k$ ] > ] \\
       arg-st & \emp{\avm{\tag{B}}} < X > $\oplus$ \tag{C} < > ] > \\
    output & < \emp{F$_{NPL}$(\avm{\1})} , [\type*{\emp{word}}
       syn & \3 [ head & [ agr & \emp{[num & pl]} ] \\
                  val & [ spr & \emp{\avm{\tag{B}}} \\ comps & \tag{C} ] ] \\
       sem & \2 \\
       arg-st & \tag{B} $\oplus$ \tag{C} ] > ]}}
\end{center}

\myslide{Three Subtypes of \textit{l-rule}}
\begin{center}
\begin{small}
  \begin{forest}
    cat, for tree={font=\itshape, l sep=3mm}
    [{l-rule} [{i-rule}] [{d-rule}] [{pi-rule}]]
  \end{forest}

  \medskip
  \val{l-rule} :
  \avm{[ input & \type{l-sequence} < X , [ sem & / \2 ] > \\
         output & \type{l-sequence} < Y , [ sem & / \2 ] > ]}

  \medskip
  \val{i-rule} :
  \avm{[ input & < X , [\type*{lexeme}
                         syn & \3 \\
                         arg-st & \tag{A} ] > \\
         output & < Y , [\type*{word}
                          syn & \3 \\
                          arg-st & \tag{A} ] > ]}
  \qquad
  \val{d-rule} :
  \avm{[ input & < X , [\type*{lexeme}
                         syn & / \3 ] > \\
         output & < Y , [\type*{lexeme}
                          syn & / \3 ] > ]}
\end{small}
\end{center}

\myslide{\textit{cats}: The Lexical Sequence}
\begin{center}
  \avm{< \textnormal{cats} , [\type*{word}
    syn & [ head & [\type*{noun} agr & \type{3pl} ] \\
            val & [ spr & \tag{B} < DP[ count & + \\ agr & \7 ] > \\
                    comps & < > ] ] \\
    sem & [ mode & ref \\
            index & $k$ \\
            restr & < [ reln & cat \\ inst & $k$ ] > ] \\
    arg-st & \tag{B} ] >}
\end{center}

\myslide{Practicalities --- Writing Lexical Rules}
\begin{itemize}
\item Determine the type of the LR.
\item Determine the class of possible inputs.
\item Determine what should change.
  \begin{small}
  \begin{itemize}
  \item If \ft{input} and \ft{output} values are identified (by default or otherwise) and
    only \ft{output} value is mentioned, then\ldots
    \\ \emp{information is added.}
    \\[1ex] (Lexical sequences incompatible with that value are not possible inputs)
  \item If \ft{input} and \ft{output} values are identified by default, but different values
    are given on the \ft{input} and \ft{output} of the rule, then\ldots
    \\ \emp{information is changed.}
  \item If \ft{input} and \ft{output} values are identified by an inviolable constraint, but
    different values are given on the \ft{input} and \ft{output} of the rule, then\ldots
    \\ \emp{there is no well-formed output}
  \end{itemize}
  \end{small}
\end{itemize}

\myslide{Constant Lexemes}
\begin{itemize}
\item What kinds of words are constant lexemes
  in our grammar?
\item Why do we need a rule for these words?
\item What would be an alternative analysis?
\end{itemize}

\myslide{Constant Lexeme LR}
\begin{center}
  \avm{[\type*{i-rule}
    input & < \1 , \type{const-lxm} > \\
    output & [ first & \1 ] ]}
\end{center}
\begin{itemize}
\item What keeps this from applying to, say, verb
  lexemes?
\item Why is this an \val{i-rule}?
\end{itemize}

\myslide{\ft{arg-st} \& ARP}
\begin{itemize}
\item Given the ARP, what do we need to specify
  about the valence properties of words?
\item Why isn't the ARP a constraint on the type
  \val{lexeme}?
\end{itemize}

\myslide{The Feature \ft{form}}
\begin{itemize}
\item Different inflected forms of verbs show
  up in different syntactic environments.
  \\ Examples?
\item These different forms are syntactically
  distinguished by the feature \ft{form}, as
  assigned by lexical rules.
\item \ft{form} is also useful in our analyses of
  coordination and PP selection.
\end{itemize}

\myslide{What Rules These Out?}
\begin{exe}
\ix \bad{Kim eat pizza.}
\ix \bad{Kim seems to eats pizza.}
\ix \bad{Dana helped Leslie pack and moved.}
\ix \bad{Kim relies for Sandy.}
\ix \bad{Dana walked and Kim.}
\end{exe}

\myslide{Overview}
\begin{itemize}
\item How lexical rules fit in
\item Three types of lexical rules, constraints
\item Example: Plural noun lexical rule
\item Advice on writing lexical rules
\item Constant lexemes
\item \ft{arg-st} \& ARP
\item The feature \ft{form}
\end{itemize}


% \myslide{P1: \eng{'s} and the SHAC}
% \MyLogo{Based on  Chapter 8 Problem 1, Sag, Wasow and Bender (2003)}

% The name `Specifier-Head Agreement Constraint' suggests that heads
% always agree with their specifiers.  Examples like \eng{Pat's parents}
% and \eng{the children's game} look like counterexamples:  in both
% cases, the possessive NP in the DP that functions as the specifier of
% the noun differs in number from that noun.

% Explain why these are not really counterexamples, given our
% formulation of SHAC as a type constraint, together with the analysis
% of possessives developed in Problem 4 of Chapter 6.  {\sl [Hint: The
%   fact that \eng{'s} is the head of the DP is crucial.]}

% \myslide{P2: Plural and Mass NPs Without Specifiers}
% \MyLogo{Based on  Chapter 8, Problem 2, Sag, Wasow and Bender (2003)}

% There is a problem with our treatment of  common nouns.  
% The type  \val{cn-lxm} requires common nouns to have nonempty SPR lists, and this
% requirement is preserved in the  Plural Noun Lexical Rule.
% Similarly, the type \val{massn-lxm} inherits the constraint on the SPR,
% and this constraint is preserved when these nouns undergo the inflectional
% rules.
% This treatment makes the wrong predictions:
% %But
% specifiers are optional for plural nouns and mass nouns.
% \begin{itemize}
% \item[A.] 
% %Give at least three examples showing the \index{optionality of specifier} optionality of 
% %specifiers for \index{noun (N)!plural noun} plural and \index{noun (N)!mass noun} mass nouns.
% Give examples showing, for one plural noun 
% and one \index{noun (N)!mass noun} mass noun, that
% the specifier is optional (i.e.\ permitted but not obligatory).
% \end{itemize}

% \newpage
% Two \index{noun phrase (NP)!without determiner|(} obvious approaches to this problem are the following:
% %\pagebreak
% \begin{itemize}
% \item[(i)] allow empty SPR lists in the lexical entries for plural
% and mass nouns; or 
% \item[(ii)] introduce a new  \index{phrase structure!rule (PS rule)}grammar rule to account for NPs with
% plural or mass heads and no specifiers.
% \end{itemize}

% \noindent
% Alternative (i) would involve modifying the Plural Noun Lexical Rule,
% %as well as introducing a new subtype of \val{cn-lxm} for mass nouns.
% as well as the type \val{massn-lxm} to make the first member of
% the ARG-ST list optional.\footnote{This would require making the
% constraint on the ARG-ST of \val{cn-lxm} defeasible.}
% \index{noun phrase (NP)!without determiner|)}  

% The rule in alternative (ii) is analogous to the \index{Imperative
% Rule} Imperative Rule given in Chapter 7, in that it would have only
% one constituent on the right hand side, and its function would be to
% license a constituent without a specifier, although its daughter has a
% nonempty SPR list.

% It turns out that alternative (i) makes incorrect predictions about
% prenominal modifiers (see Problem~1 of Chapter~5).  We want adjectives
% like \eng{cute} to modify plural nouns even when they don't have
% specifiers:

% \begin{itemize}
% \item[(iii)] \eng{Cute puppies make people happy.}
% \end{itemize}

% \noindent
% Under alternative (i), in order to generate (iii), we would have
% to allow adjectives like \eng{cute} to modify NPs (i.e.\ expressions
% that are \mbox{[SPR $\langle$  $\rangle$]}).  If we do that, however, we have no
% way to block (iv):\footnote{There are also technical problems with
% making alternative (i) work with the ARP.}

% \begin{itemize}
% \item[(iv)] \bad{Cute the puppies make people happy.}
% \end{itemize}

% \newpage
% Alternative (ii), on the other hand, would allow \eng{cute}
% to always modify a \index{NOM}NOM \mbox{([SPR $\langle$ DP $\rangle$])}
% constituent.  A NOM, modified or otherwise, could either be
% the daughter of the non-branching rule, or the head daughter of
% the \index{Head Specifier Rule@Head-Specifier Rule}Head-Specifier Rule.


% \begin{itemize}
% %\item[B.]  The ARP does not manage to sufficiently constrain what
% %alternative (i) would allow.  Explain why.\smallskip
% %
% %\noindent
% %[{\sl Hint: Assume that alternative
% %(i) involves making the \index{determiner (D)} determiner an optional member of the noun's
% %ARG-ST. Consider all the ways that an \index{ARGUMENT STRUCTURE (ARG ST)@ARGUMENT-STRUCTURE (ARG-ST)} ARG-ST element could be realized
% %on the SPR or COMPS list.  You may also want to construct examples of
% %common nouns with complements but no specifiers.}]

% \item[B.]  Formulate the rule required for alternative (ii).  \smallskip
% %Be as precise as you can.\smallskip%\hfill\break

% \noindent
% {[{\sl Hint:  The trickiest part is formulating the rule so that it
% applies to both plural \index{noun (N)!count noun}count nouns and \index{noun (N)!mass noun}mass nouns, while not applying
% to singular count nouns.  You will need to include a disjunction in 
% the rule.  The 
% %\index{SPR}
% SPR list of the \index{head daughter (H)}head daughter is a good place to state
% it, since the three types of nouns differ in the requirements they
% place on their specifiers.}]}
% \index{noun (N)!singular noun}
% \index{specifier!optional|)}

% %disjunction of plural or mass nouns.  The place to do this is in the
% %SPR list of the feature structure description 
% %on the right hand side of the rule.}]
% \end{itemize}


\section{Problem}

\myslide{P1: Arguments in Japanese}
\MyLogo{Based on \citet[Chapter 8, Problem 6]{sag:wasow:bender:2003}}
As noted in Chapter 2, Japanese word order differs from English in a
number of ways, including the fact that it is a `Subject-Object-Verb'
(SOV) language.  The next two slides give some relevant examples. In
the glosses, `{\sc nom}', `{\sc acc}', and `{\sc dat}' stand for
nominative, accusative, and dative case, respectively.  (Note that
Japanese has one more case --- dative --- than English does.  This
doesn't have any important effects on the analysis; it merely requires
that we posit one more possible value of CASE for Japanese than for
English.)\footnote{The examples marked with `*' here are unacceptable
  with the indicated meanings.  Some of these might be well-formed
  with some other meaning of no direct relevance; others might be
  well-formed with special intonation that we will ignore for present
  purposes.}

\myslide{P1: Arguments in Japanese --- Data}
\begin{exe}
  \ex\label{jp:first} %一人 の 男 が その 本 を 読んだ 。
  \gll Hitorino  otoko-ga sono  hon-o yonda. \\
 one  man-{\sc nom} that  book-{\sc acc} read.{\sc pst} \\
\trans {`One man read that book.'}

[cf.\ *Yonda hitorino otoko-ga sono hon-o. \\ {*}Hitorino otoko-ga yonda
sono hon-o. \\ {*}Otoko-ga hitorino sono hon-o yonda. \\ {*}Hitorino
otoko-ga hon-o sono yonda. \\ {*}Hitorino otoko-ni/-o sono hon-o
yonda. \\ {*}Hitorino otoko-ga sono hon-ga/-ni yonda.]
\end{exe}

\myslide{P1: Arguments in Japanese --- Data (continued)}
\begin{exe}
\ex 
\gll  Hanako-ga  hon-o  yonda. \\
 Hanako-{\sc nom}  book-{\sc acc}  read.{\sc pst} \\
\trans {`Hanako read the book(s).'}

[cf.\ *Yonda Hanako-ga hon-o. \\ {*}Hanako-ga yonda hon-o. \\
{*}Hanako-ni/-o hon-o yonda. \\ {*}Hanako-ga hon-ni/-ga yonda.]
\end{exe}

\myslide{P1: Arguments in Japanese --- More Data}
\begin{exe}
\ex
\gll   Sensei-ga  Taroo-ni  sono   hon-o  ageta. \\
 teacher-{\sc nom} Taroo-{\sc dat} that   book-{\sc acc} give.{\sc pst} \\
 \trans {`The teacher(s) gave that book to Taroo.'}

[cf.\ *Ageta sensei-ga Taroo-ni sono hon-o. \\ {*}Sensei-ga ageta Taroo-ni
sono hon-o. \\ {*}Sensei-ga Taroo-ni ageta sono hon-o. \\ {*}Sensei-o/-ni
Taroo-ni sono hon-o ageta. \\ {*}Sensei-ga Taroo-ga/-o sono hon-o ageta.\\
{*}Sensei-ga Taroo-ni sono hon-ga/-ni ageta.]

\ex\label{jp:last}
\gll Hanako-ga   kita. \\
 Hanako-{\sc nom}   arrive.{\sc pst} \\
\trans {`Hanako arrived.'}

[cf.\ *Kita Hanako-ga.]
\end{exe}

\myslide{P1: Arguments in Japanese --- Tasks}
As the contrasting ungrammatical examples show, the verb must appear
in final position in Japanese. In addition, we see that verbs select
for NPs of a particular case, much as in English. In the following
tasks, assume that the nouns and verbs of Japanese are inflected
words, derived by lexical rule from the appropriate lexemes.

\begin{itemize}\addtolength{\itemsep}{1ex}
\item[A.] Write Head-Specifier and Head-Complement Rules for
Japanese that account for the data illustrated here.  How are
they different (if at all) from the Head-Specifier and Head-Complement 
Rules for English?
\end{itemize}

\myslide{P1: Arguments in Japanese --- Tasks (continued)}
\begin{itemize}\addtolength{\itemsep}{1ex}
\item[B.] Give the lexical entry for each of the verbs
illustrated in (\ref{jp:first})--(\ref{jp:last}). 

[{\sl Make sure your entries interact with the rules you
formulated in part (A) to account for the above data.  The data given
permit you to specify only some features; leave others unspecified.
Assume that there is a Past-Tense Verb Lexical Rule (an \val{i-rule})
that relates your lexical entries to the words shown in (\ref{jp:first})--(\ref{jp:last}).  We
have not provided a hierarchy of lexeme types for Japanese.  You may
either give all relevant constraints directly on the lexical entries,
or posit and use subtypes of \val{lexeme}.  In the latter case, you
must also provide those types.}]

\item[C.] Give the lexical entries for the nouns \eng{Taroo} and
\eng{hon}. [{\sl See notes on part (B).}]

\item[D.] Formulate the lexical rule for deriving the inflected
forms ending in \eng{-o} from the nominal lexemes.
\end{itemize}


\myslide{Acknowledgments and References}

\begin{itemize}
\item Course design and slides borrow heavily from Emily Bender's course:
\textit{Linguistics 566: Introduction to Syntax for Computational Linguistics}
\\ \url{http://courses.washington.edu/ling566}
\item The slides on lexical types and lexical rules are adapted from
  Emily Bender's, with her permission; the feature structures follow
  \citet{sag:wasow:bender:2003}, \copyright{} 2003 CSLI Publications
\item This lecture complements \citet[Chapter 8]{sag:wasow:bender:2003}
\end{itemize}


%\input{lec-05-questions}

\bibliographystyle{aclnat}
\bibliography{abb,mtg,nlp,ling}
\end{document}

%%% Local Variables: 
%%% coding: utf-8
%%% mode: latex
%%% TeX-PDF-mode: t
%%% TeX-engine: xetex
%%% End: 
